quotient of @@M0@@ by the @@M1@@-st factor. In this system of morphisms, the canonical quotients therefore correspond by inverse image. The inductive system thus defined will be denoted by @@M2@@.

Its @@M3@@ is defined by the formula:
\sgaStableTarget{sga6.e06.eq.4-9-1}@@M4@@

Let @@M5@@ be the canonical sheaves on @@M6@@ (4.5), and @@M7@@ their class in @@M8@@. Put @@M9@@. As @@M10@@ varies, the @@M11@@ correspond under the transition morphisms, and therefore define elements, still denoted by @@M12@@, of @@M13@@.

\sgaStableTarget{sga6.e06.prop.4-10}\textbf{Proposition 4.10.} \uline{Let} @@M14@@ \uline{be a quasi}-\uline{compact scheme}. \uline{Then}, \uline{with the notation of} 4.9, \uline{there is an isomorphism}:
@@M15@@

By 4.8, @@M16@@ is the @@M17@@-algebra generated by the @@M18@@, @@M19@@, and the @@M20@@, @@M21@@, subject to the relations deduced from:
@@M22@@
or again:
@@M23@@

This relation expresses the @@M24@@ in terms of the @@M25@@, and gives a series of isobaric relations among the @@M26@@ (@@M27@@ being assumed to have weight @@M28@@) by writing that the coefficients of @@M29@@ are zero in degree @@M30@@; moreover, these relations generate the ideal of relations among the @@M31@@. We may therefore write:
